\Askhsh{34505}{2}
\begin{ylh}{1}
    \item 3.1. Είδη και στοιχεία τριγώνων
    \item 3.6. Κριτήρια ισότητας ορθογώνιων τριγώνων
    \item 5.2. Παραλληλόγραμμα
    \item 5.4. Ρόμβος.
\end{ylh}
% ===================================================================
%  Αρχείο: 34505.tex (Fragment)
%  Αυτόματη μετατροπή μέσω doc2latex.py
% ===================================================================

ΘΕΜΑ 2

Σε ημικύκλιο διαμέτρου ΑΒ προεκτείνουμε την ΑΒ προς το μέρος του Α και παίρνουμε ένα σημείο Γ. Θεωρούμε Ε ένα σημείο του ημικυκλίου και έστω Δ το σημείο τομής του τμήματος ΓΕ με το ημικύκλιο. Αν το τμήμα ΓΔ είναι ίσο με το ΟΒ και η γωνία $Β\widehat{\Gamma}Ε\ \varepsilon ί\nu\alpha\iota\ 15{^\circ}$, τότε

\begin{alist}
\item να αποδείξετε ότι $O\widehat{\Delta}E = 30^{\circ}$. \monades{13}
\item να υπολογίσετε τη γωνία $E\widehat{O}B = x$. \monades{12}

\begin{center}
\begin{tikzpicture}[scale=1]
\tkzDefPoint(0,0){O}
\tkzDefPoint(3,0){B}
\tkzDefPoint(-3,0){A}
\tkzDrawArc[l3](O,B)(A)
\draw[l3] (A)--(B);

\tkzDefPointBy[homothety=center O ratio 2](A) \tkzGetPoint{C}
\tkzDefPoint(45:3){E} % Just an example E
% Actually ODE=30, angle C=15.
% CD=OB=r.
% Let C = (-6,0). O=(0,0). E is such that angle ECO=15.
\tkzDefPoint(-6,0){C}
\tkzDefPointBy[rotation=center C angle 15](O) \tkzGetPoint{line_point}
\tkzInterLC(C,line_point)(O,B) \tkzGetPoints{E}{D} 

\draw[l3] (C)--(E);
\draw[l3] (O)--(D);
\draw[l3] (O)--(E);
\draw[l3] (A)--(C);

\tkzDrawPoints(O,A,B,C,D,E)
\tkzLabelPoint[below](O){$O$}
\tkzLabelPoint[below](A){$A$}
\tkzLabelPoint[below](B){$B$}
\tkzLabelPoint[below](C){$\Gamma$}
\tkzLabelPoint[above left](D){$\Delta$}
\tkzLabelPoint[above right](E){$E$}

\tkzMarkAngle[size=1.2,mark=none](O,C,E)
\tkzLabelAngle[pos=1.5](O,C,E){$15^\circ$}
\tkzMarkAngle[size=0.6,mark=none](B,O,E)
\tkzLabelAngle[pos=0.9](B,O,E){$x$}
\tkzMarkSegments[mark=s|](C,D O,B O,A O,D O,E)
\end{tikzpicture}
\end{center}
\end{alist}
